\subsection*{Associativity}
\label{subsec:associativity}

\begin{definitionbox}{Definition (Associativity)}
\begin{definition}[Associativity]
\label{def:algebraic-associativity}
Let $\star$ be a binary operation on $G$.
The operation $\star$ is \textbf{associative} if
\[
(a \star b) \star c = a \star (b \star c)
\quad \text{for all } a,b,c \in G.
\]
\end{definition}
\end{definitionbox}

\begin{remark*}[Standard quantified statement]
\[
\operatorname{Associative}(\star,G)
\;\colonequiv\;
\forall a,b,c \in G,\;
(a \star b) \star c = a \star (b \star c).
\]
\end{remark*}

\begin{remark*}[Definition predicate reading]
$\star$ is associative on $G$ exactly when any threefold product has the same
value under either parenthesization.
\end{remark*}

\begin{remark*}[Negated quantified statement]
\[
\neg\operatorname{Associative}(\star,G)
\;\colonequiv\;
\exists a,b,c\in G,\;
(a\star b)\star c\neq a\star(b\star c).
\]
\end{remark*}

\begin{remark*}[Interpretation]
Associativity says that parenthesization does not affect the result of a
chain of operations.
\end{remark*}

\begin{dependencies}
\begin{itemize}
  \item \hyperref[def:binary-operation]{Binary Operation}
\end{itemize}
\end{dependencies}
